% Part: lambda-calculus
% Chapter: lambda-definability
% Section: truth-values

\documentclass[../../../include/open-logic-section]{subfiles}

\begin{document}

\olfileid{lam}{ldf}{tvr}
\olsection{सत्यतामूल्ये आणि संबंध}

शुद्ध लॅम्डा कलनात सत्यतामूल्यांचे पुढीलप्रमाणे
संकेतन करता येते:
\begin{align*}
  \fn{true} & \ident \lambd[x][\lambd[y][x]]\\
  \fn{false} & \ident \lambd[x][\lambd[y][y]]
\end{align*}

सत्यतामूल्ये \emph{निवडफलने} म्हणून दर्शवली जातात:
दोन फलसाधक स्वीकारून त्यांपैकी एक परत करणारी
फलने. सत्यतामूल्य $\fn{true}$ पहिला फलसाधक
निवडते आणि $\fn{false}$ दुसरा. उदाहरणार्थ,
$\fn{true}\, M N$ चे नेहमी $M$ मध्ये, तर
$\fn{false}\, M N$ चे नेहमी~$N$ मध्ये न्यूनीकरण होते.

\begin{defn}
संबंध $R \subseteq \Nat^k$ हा !!{lambda definable}
आहे असे म्हणतो, जर आणि फक्त जर असे बंद पद~$R$ असते की
प्रत्येक नैसर्गिक निविष्टांसाठी पुढील अटी पूर्ण होतात:
\begin{align*}
  R\, \num{n_1} \dots \num{n_k} & \bred \fn{true}
  \intertext{जेव्हा $R(n_1, \dots, n_k)$ सत्य असेल, आणि}
  R\, \num{n_1} \dots \num{n_k} & \bred \fn{false}
\end{align*}
अन्यथा.
\end{defn}

उदाहरणार्थ, $0$ साठी खरा आणि $0$ व्यतिरिक्त इतर
कोणत्याही संख्येसाठी खरा नसलेला संबंध
$\fn{IsZero} = \{0\}$ पुढील पदाने
!!{lambda definable} आहे:
\[
  \fn{IsZero} \ident \lambd[n][n (\lambd[x][\fn{false}])\, \fn{true}].
\]
हे कसे कार्य करते? चर्च अंक हे पुनरावर्तक आहेत:
पहिला फलसाधक दुसऱ्यावर $n$~वेळा लागू करणारी
फलने. म्हणून सुरुवातीचे मूल्य $\fn{true}$ घेऊ;
पुनरावर्तनाच्या प्रत्येक पायरीत आधीच्या पायरीचा
निकाल काहीही असला, तरी $\fn{false}$ परत करू.
ही पायरी सुरुवातीच्या मूल्यावर $n$ वेळा लागू होते.
म्हणून ती एकदाही लागू झाली नाही, म्हणजे $n = 0$
असेल, तेव्हाच निकाल $\fn{true}$ असतो.

सत्यतामूल्यांच्या या संकेतनावरून आणखी सत्यता-फलने
परिभाषित करता येतात. नकरण आणि संधियोग यांची
दोन निरूपणे अशी:
\begin{align*}
  \fn{Not} & \ident \lambd[x][x\, \fn{false}\, \fn{true}]\\
  \fn{And} & \ident \lambd[x][\lambd[y][xy \,\fn{false}]]
\end{align*}
``$\fn{Not}$'' हे एक फलसाधक स्वीकारते:
तो $\fn{false}$ असेल, तर $\fn{true}$ परत करते;
आणि तो~$\fn{true}$ असेल, तर $\fn{false}$.
``$\fn{And}$'' हे दोन सत्यतामूल्ये स्वीकारते आणि
निकाल $\fn{true}$ असायला हवा, जर आणि फक्त जर
दोन्ही निविष्ट सत्यतामूल्ये~$\fn{true}$ असतील. वर सांगितल्याप्रमाणे सत्यतामूल्ये
निवडफलने आहेत. म्हणून सत्यतामूल्य $x$ ला दोन
फलसाधक लागू केल्यास, $x$ हे $\fn{true}$ असेल
तर पहिला फलसाधक आणि अन्यथा दुसरा मिळतो.
आता $\fn{And}$ आपले दोन फलसाधक $x$ आणि $y$
घेते व $x$ ला अनुक्रमे $y$ आणि $\fn{false}$
हे देते. $x$ सत्यतामूल्य आहे असे गृहीत धरल्यास,
$x$ हे $\fn{true}$ असेल तर निकाल~$y$,
आणि $x$ हे $\fn{false}$ असेल तर निकाल
$\fn{false}$ होतो; हाच अपेक्षित परिणाम आहे.

येथे $\fn{And}$ ला फलसाधक म्हणून फक्त
सत्यतामूल्येच दिली जातात असे आपण गृहीत धरतो.
इतर पदे दिल्यास निकाल—सामान्य रूप अस्तित्वात
असेल तर ते—सत्यतामूल्य असेलच असे नाही.

\begin{prob}
  सत्यतामूल्यांचे $\lambd$-पदांद्वारे केलेले संकेतन
  वापरून समावेशक आणि अपवर्जक विकल्पयोग दर्शवणारी
  $\fn{Or}$ आणि $\fn{Xor}$ ही फलने परिभाषित करा.
\end{prob}

\end{document}
